\begin{table}[H]
\centering
\caption{Review of source vs. story approaches.}
\label{table:table1}
\begin{tabularx}{\textwidth}{ | l | X | X | }
\toprule
        &  \textbf{Story} &  \textbf{Source}           \\
\midrule
\textbf{Reliability} &  
+ high precision in story accuracy ratings \newline
– bias in story selection reflected in biased estimates 
&
+ agreement across different lists for the US context (Lasser at al., 2022) \newline
+ source ratings can still be unbiased at the macroscale and represent variance in story accuracy
\\
 \hline
\textbf{Validity}    &  
– not generalizable across topics and context \newline
– topics change fast: timing most likely causes lags in the coverage of novel content \newline
– often binary classifications (true vs false news) \newline
+ high construct validity
 & 
+ generalizable across topics
- not across countries: journalistic norms depend on media system \newline
+ can take biases and tendencies into consideration \newline
– low construct validity \newline
– excludes misinformation spread by individual social media users 
\\
 \hline
\textbf{Feasibility} &  
– scalable up to a certain degree: needs automated approaches or crowdsourced initiatives like Birdwatch \newline
– needs experts for all topics, especially if they change over time \newline
– information available for fact-checked stories
& 
+ scalable at very high macrolevel \newline
+ expert vs. crowdsourced ratings are similar (Pennycook \& Rand, 2019) \newline
+ lists with source ratings more widely available at the publisher- than article-level \newline
– few open-source lists
\\
\bottomrule
\end{tabularx}
\end{table}